\iffalse
% \section{Дополнительные структуры: соглашения}
% Все листинги самостоятельны; снаружи заданы \texttt{bits/stdc++.h},
% \texttt{using namespace std; using ll=long long; using db=double; ll INF=1e9;}.
% Глобальный \texttt{INF} не используется как бесконечность для \texttt{ll}.
% Индексы начинаются с нуля, отрезки полуоткрытые. Вход удовлетворяет
% описанным предусловиям. Размеры массивов, индексы узлов и их суммы помещаются
% в \texttt{int}; память должна вмещать создаваемые узлы. Все суммы,
% разности, произведения и промежуточные результаты помещаются в \texttt{ll};
% расширенная целочисленная арифметика не используется.
% Деревья --- связные неориентированные простые графы; если не сказано иначе,
% они непусты и укоренены в вершине 0. Рекурсивным реализациям нужен стек
% глубины $O(n)$ (у случайных treap --- ожидаемо $O(\log n)$).
% Листинги написаны самостоятельно, а не скопированы из Algorithmica.
% Тематические материалы Algorithmica по DSU и деревьям Фенвика:
% \url{https://ru.algorithmica.org/cs/set-structures/dsu/},
% \url{https://ru.algorithmica.org/cs/range-queries/fenwick/}.
% Эти ссылки не являются источниками всех приведённых ниже алгоритмов.
\fi
\section{Система непересекающихся множеств с откатом}
Изначально \texttt{parent[i]=i}, \texttt{size[i]=1}, история пуста; снимок --- её размер.
Объединение возвращает признак успеха, откат восстанавливает снимок только текущей ветви истории.
Без сжатия путей поиск и объединение стоят $O(\log n)$, откат --- $O(1)$ на отменённое объединение; память $O(n)$.
\begin{lstlisting}[name=rollback_dsu]
int rollback_root(const vector<int>& parent, int vertex) {
    while (parent[vertex] != vertex) vertex = parent[vertex];
    return vertex;
}
bool rollback_unite(vector<int>& parent, vector<int>& size,
        vector<array<int, 3>>& history, int first, int second) {
    first = rollback_root(parent, first);
    second = rollback_root(parent, second);
    if (first == second) return false;
    if (size[first] < size[second]) swap(first, second);
    history.push_back({second, first, size[first]});
    parent[second] = first;
    size[first] += size[second];
    return true;
}
void rollback(vector<int>& parent, vector<int>& size,
        vector<array<int, 3>>& history, int snapshot) {
    assert(0 <= snapshot && snapshot <= int(history.size()));
    while (int(history.size()) > snapshot) {
        auto [child, root, oldSize] = history.back();
        history.pop_back();
        parent[child] = child;
        size[root] = oldSize;
    }
}
\end{lstlisting}
% Позволяет вернуться к состоянию до группы объединений, например в офлайн
% динамической связности. Изначально \texttt{parent[i]=i}, \texttt{size[i]=1},
% история пуста. Снимок --- размер истории; откатывать можно только к снимку
% текущей ветви истории. Сжатия путей нет. Поиск и объединение стоят
% $O(\log n)$, откат --- $O(1)$ на отменённое успешное объединение;
% память $O(n)$. Неудачное объединение историю не меняет.
%
% Историю можно дополнить агрегатами компонент и числом компонент. Для
% динамической связности интервалы жизни рёбер раскладывают по дереву времени.

\section{Система непересекающихся множеств с чётностью}
Ограничение $(u,v,p)$ задаёт $color[u]\mathbin{\mathrm{xor}}color[v]=p$, $p\in\{0,1\}$; вершины нумеруются с нуля.
Возвращает совместность каждого префикса ограничений; после противоречия все ответы ложны, для двудольности берут $p=1$.
Время $O((n+q)\alpha(n))$, память $O(n+q)$.
\begin{lstlisting}[name=parity_dsu]
vector<bool> parity_dsu(int n, const vector<array<int, 3>>& edges) {
    vector<int> parent(n), size(n, 1), parity(n);
    iota(parent.begin(), parent.end(), 0);
    auto find = [&](auto&& self, int vertex) -> int {
        if (parent[vertex] == vertex) return vertex;
        int oldParent = parent[vertex];
        parent[vertex] = self(self, oldParent);
        parity[vertex] ^= parity[oldParent];
        return parent[vertex];
    };
    vector<bool> answer;
    bool valid = true;
    for (auto [first, second, required] : edges) {
        int a = find(find, first), b = find(find, second);
        int difference = parity[first] ^ parity[second] ^ required;
        if (a == b) valid = valid && difference == 0;
        else {
            if (size[a] < size[b]) swap(a, b);
            parent[b] = a;
            parity[b] = difference;
            size[a] += size[b];
        }
        answer.push_back(valid);
    }
    return answer;
}
\end{lstlisting}
% Пакет ограничений $(u,v,p)$ задаёт $color[u]\mathbin{\mathrm{xor}}color[v]=p$,
% где $p\in\{0,1\}$. После каждого ограничения возвращается признак
% совместности всего префикса; после противоречия все ответы ложны.
% Для проверки двудольности каждое ребро имеет $p=1$.
% Время $O((n+q)\alpha(n))$, память $O(n+q)$.
%
% Можно хранить корректность отдельно для каждой компоненты или заменить
% чётность потенциалами по модулю. Для отката нужно отказаться от сжатия путей.

\section{Два дерева Фенвика: прибавление и сумма на отрезке}
Массив изначально содержит $n$ нулей; запрос $(type,left,right,delta)$ типа 0 прибавляет \texttt{delta} на $[left,right)$, типа 1 возвращает сумму этого отрезка.
Ответы выдаются только для типа 1; требуется $0\le n<\texttt{INT\_MAX}$, все суммы, отрицание \texttt{delta} и произведения на индексы помещаются в \texttt{ll}.
Время $O((n+q)\log(n+1))$, память $O(n+q)$.
\begin{lstlisting}[name=fenwick_range]
vector<ll> fenwick_range(int n,
        const vector<tuple<int, int, int, ll>>& queries) {
    vector<ll> slope(n + 1), offset(n + 1), answer;
    auto add = [&](vector<ll>& tree, int index, ll delta) {
        for (++index; index <= n;) {
            tree[index] += delta;
            int step = index & -index;
            if (step > n - index) break;
            index += step;
        }
    };
    auto prefix = [&](int right) {
        ll a = 0, b = 0;
        for (int index = right; index > 0; index -= index & -index) {
            a += slope[index]; b += offset[index];
        }
        return a * right - b;
    };
    for (auto [type, left, right, delta] : queries) {
        if (type == 0) {
            add(slope, left, delta); add(slope, right, -delta);
            add(offset, left, delta * left);
            add(offset, right, -delta * right);
        } else answer.push_back(prefix(right) - prefix(left));
    }
    return answer;
}
\end{lstlisting}
% Исходный массив из $n\ge0$ нулей. Запрос $(type,left,right,delta)$:
% тип 0 прибавляет \texttt{delta} на $[left,right)$, тип 1 возвращает сумму
% (\texttt{delta} игнорируется). Ответы только для типа 1.
% Время $O((n+q)\log(n+1))$, память $O(n+q)$; $n<\texttt{INT\_MAX}$.
% В \texttt{ll} должны помещаться также отрицание \texttt{delta},
% произведения разностей на индексы и суммы в обоих деревьях.
%
% Начальный массив можно построить по разностям за линейное время. Для
% присваиваний вместо прибавлений удобнее дерево с отложенными операциями.

\section{Поиск порядковой статистики в дереве Фенвика}
Для массива неотрицательных частот и рангов $k\ge1$ возвращает индексы первых префиксов с суммой $\ge k$; если общей суммы не хватает, ответ $n$.
Пустой массив допустим; $n<\texttt{INT\_MAX}$, частоты и их сумма помещаются в \texttt{ll}.
Построение $O(n)$, запрос $O(\log(n+1))$, память $O(n+q)$.
\begin{lstlisting}[name=fenwick_kth]
vector<int> fenwick_kth(const vector<ll>& frequencies,
        const vector<ll>& ranks) {
    int n = frequencies.size();
    vector<ll> tree(n + 1);
    for (int i = 1; i <= n; ++i) {
        tree[i] += frequencies[i - 1];
        int step = i & -i;
        if (step <= n - i) tree[i + step] += tree[i];
    }
    int highest = 1;
    while (highest <= n / 2) highest *= 2;
    vector<int> answer;
    for (ll rank : ranks) {
        assert(rank >= 1);
        int index = 0;
        for (int step = highest; step > 0; step /= 2) {
            if (step <= n - index && tree[index + step] < rank) {
                index += step;
                rank -= tree[index];
            }
        }
        answer.push_back(index);
    }
    return answer;
}
\end{lstlisting}
% Для неизменяемых неотрицательных частот строит дерево и для каждого
% $k\ge1$ возвращает индекс первого префикса с суммой не меньше $k$.
% Если $k$ больше общей суммы, возвращает $n$; пустой массив допустим.
% Построение $O(n)$, запрос $O(\log(n+1))$, память $O(n+q)$.
% Все частоты и их общая сумма помещаются в \texttt{ll},
% $n<\texttt{INT\_MAX}$. Отрицательные частоты запрещены.
%
% Для онлайн-частот построение и поиск можно вынести в отдельные функции с
% общим внешним массивом; точечные изменения обязаны сохранять неотрицательность.

\section{Разреженная таблица минимумов}
Для неизменяемого массива возвращает минимумы на непустых полуинтервалах $[left,right)$; пустой массив допустим только без запросов.
Построение и память $O(n\log(n+1))$, запрос $O(1)$, ответы занимают $O(q)$ памяти.
\begin{lstlisting}[name=sparse_table]
vector<ll> sparse_table(const vector<ll>& values,
        const vector<pair<int, int>>& queries) {
    int n = values.size();
    vector<int> logarithm(n + 1);
    for (int i = 2; i <= n; ++i) logarithm[i] = logarithm[i / 2] + 1;
    vector<vector<ll>> table(logarithm[n] + 1, values);
    for (int level = 1; level <= logarithm[n]; ++level) {
        int length = 1 << level, half = length / 2;
        for (int i = 0; i <= n - length; ++i)
            table[level][i] = min(table[level - 1][i],
                                  table[level - 1][i + half]);
    }
    vector<ll> answer;
    for (auto [left, right] : queries) {
        int level = logarithm[right - left];
        answer.push_back(min(table[level][left],
                             table[level][right - (1 << level)]));
    }
    return answer;
}
\end{lstlisting}
% Минимум неизменяемого массива на непустых $[left,right)$ за $O(1)$.
% Построение и память $O(n\log(n+1))$, ответы занимают ещё $O(q)$.
% Пустой массив разрешён только без запросов.
%
% Минимум заменяется любой идемпотентной ассоциативной операцией, например
% максимумом или НОД. Для суммы перекрывающиеся блоки не подходят;
% можно использовать непересекающуюся разреженную таблицу.

\section{Персистентное дерево: порядковая статистика отрезка}
Для массива и запросов $(left,right,k)$ возвращает $k$-е по возрастанию значения на $[left,right)$ с учётом повторов.
Требуется $1\le k\le right-left$; при пустом массиве запросов нет, число узлов помещается в \texttt{int}.
Время $O((n+q)\log(n+1))$, память $O(n\log(n+1)+q)$.
\begin{lstlisting}[name=persistent_kth]
vector<ll> persistent_kth(const vector<ll>& values,
        const vector<array<int, 3>>& queries) {
    vector<ll> coordinates = values;
    sort(coordinates.begin(), coordinates.end());
    coordinates.erase(unique(coordinates.begin(), coordinates.end()),
                      coordinates.end());
    struct Node { int left = 0, right = 0, count = 0; };
    vector<Node> nodes(1);
    auto insert = [&](auto&& self, int old, int left, int right,
                      int position) -> int {
        int node = nodes.size();
        Node copy = nodes[old];
        nodes.push_back(copy);
        ++nodes[node].count;
        if (right - left > 1) {
            int middle = left + (right - left) / 2;
            if (position < middle) {
                int child = self(self, copy.left, left, middle, position);
                nodes[node].left = child;
            } else {
                int child = self(self, copy.right, middle, right, position);
                nodes[node].right = child;
            }
        }
        return node;
    };
    int length = coordinates.size();
    vector<int> roots(1, 0);
    for (ll value : values) {
        int position = lower_bound(coordinates.begin(), coordinates.end(), value)
                     - coordinates.begin();
        roots.push_back(insert(insert, roots.back(), 0, length, position));
    }
    vector<ll> answer;
    for (auto [left, right, rank] : queries) {
        int a = roots[left], b = roots[right], low = 0, high = length;
        while (high - low > 1) {
            int count = nodes[nodes[b].left].count - nodes[nodes[a].left].count;
            int middle = low + (high - low) / 2;
            if (rank <= count) {
                a = nodes[a].left; b = nodes[b].left; high = middle;
            } else {
                rank -= count;
                a = nodes[a].right; b = nodes[b].right; low = middle;
            }
        }
        answer.push_back(coordinates[low]);
    }
    return answer;
}
\end{lstlisting}
% Запрос $(left,right,k)$ возвращает $k$-й по возрастанию элемент
% $[left,right)$ с учётом повторов, $1\le k\le right-left$.
% Версии хранят частоты префиксов после сжатия координат.
% Время $O(n\log(n+1)+q\log(n+1))$, память $O(n\log(n+1)+q)$.
% При $n=0$ запросов нет. Счётчики ограничены $n$.
%
% В узлах можно хранить суммы для суммы $k$ наименьших. Версии по вершинам
% дерева вместе с LCA дают порядковые статистики на пути.

\section{Динамическое дерево отрезков}
Нулевой массив задан на целочисленном $[lower,upper)$; запрос $(type,left,right,delta)$ типа 0 прибавляет \texttt{delta} к позиции \texttt{left}, типа 1 возвращает сумму на $[left,right)$.
Требуется $lower<upper$; разность границ и все суммы помещаются в \texttt{ll}, число узлов --- в \texttt{int}.
Время $O(q\log(upper-lower+1))$, память $O(1+u\log(upper-lower+1)+q)$ для $u$ обновлений; запросы узлов не создают.
\begin{lstlisting}[name=dynamic_segment]
vector<ll> dynamic_segment(ll lower, ll upper,
        const vector<tuple<int, ll, ll, ll>>& queries) {
    struct Node { int left = 0, right = 0; ll sum = 0; };
    vector<Node> nodes(2);
    auto add = [&](auto&& self, int node, ll left, ll right,
                   ll position, ll delta) -> void {
        nodes[node].sum += delta;
        if (right - left == 1) return;
        ll middle = left + (right - left) / 2;
        bool goLeft = position < middle;
        int child = goLeft ? nodes[node].left : nodes[node].right;
        if (!child) {
            child = nodes.size(); nodes.push_back(Node{});
            if (goLeft) nodes[node].left = child;
            else nodes[node].right = child;
        }
        self(self, child, goLeft ? left : middle, goLeft ? middle : right,
             position, delta);
    };
    auto sum = [&](auto&& self, int node, ll left, ll right,
                   ll queryLeft, ll queryRight) -> ll {
        if (!node || queryRight <= left || right <= queryLeft) return 0;
        if (queryLeft <= left && right <= queryRight) return nodes[node].sum;
        ll middle = left + (right - left) / 2;
        return self(self, nodes[node].left, left, middle, queryLeft, queryRight)
             + self(self, nodes[node].right, middle, right, queryLeft, queryRight);
    };
    vector<ll> answer;
    for (auto [type, left, right, delta] : queries) {
        if (type == 0) add(add, 1, lower, upper, left, delta);
        else answer.push_back(sum(sum, 1, lower, upper, left, right));
    }
    return answer;
}
\end{lstlisting}
% Разреженный нулевой массив на целочисленном $[lower,upper)$.
% Запрос $(type,left,right,delta)$: тип 0 прибавляет к позиции
% \texttt{left} (\texttt{right} игнорируется), тип 1 запрашивает сумму
% $[left,right)$ (\texttt{delta} игнорируется).
% Требуется $lower<upper$, разность $upper-lower$ помещается в \texttt{ll}.
% Время $O(q\log(upper-lower+1))$, память $O(1+u\log(upper-lower+1)+q)$,
% где $u$ --- число обновлений; запросы узлов не создают.
%
% При известных заранее координатах проще сжать их. Для обновлений отрезка
% можно добавить отложенные метки, учитывая длины отсутствующих поддеревьев.

\iffalse
\section{Неявное декартово дерево с разворотом}
Для массива и последовательности полуинтервалов $[left,right)$ возвращает массив после их разворотов; пустые отрезки допустимы.
Ожидаемое время $O((n+q)\log(n+1))$, память $O(n)$; худшее время квадратичное, глубина рекурсии до $O(n)$.
Фиксированный seed не защищает от специально подобранного входа.
\begin{lstlisting}[name=implicit_treap]
vector<ll> implicit_treap(const vector<ll>& values,
        const vector<pair<int, int>>& reversals) {
    struct Node {
        ll value; int priority;
        int left = 0, right = 0, size = 1;
        bool reversed = false;
    };
    vector<Node> nodes(1, Node{0, 0});
    nodes[0].size = 0;
    mt19937 random(712367);
    auto pull = [&](int node) {
        nodes[node].size = 1 + nodes[nodes[node].left].size
                            + nodes[nodes[node].right].size;
    };
    auto push = [&](int node) {
        if (!node || !nodes[node].reversed) return;
        swap(nodes[node].left, nodes[node].right);
        for (int child : {nodes[node].left, nodes[node].right})
            if (child) nodes[child].reversed = !nodes[child].reversed;
        nodes[node].reversed = false;
    };
    auto split = [&](auto&& self, int node, int count) -> pair<int, int> {
        if (!node) return {0, 0};
        push(node);
        int leftSize = nodes[nodes[node].left].size;
        if (count <= leftSize) {
            auto [first, second] = self(self, nodes[node].left, count);
            nodes[node].left = second; pull(node);
            return {first, node};
        }
        auto [first, second] = self(self, nodes[node].right, count-leftSize-1);
        nodes[node].right = first; pull(node);
        return {node, second};
    };
    auto merge = [&](auto&& self, int first, int second) -> int {
        if (!first || !second) return first ? first : second;
        push(first); push(second);
        if (nodes[first].priority > nodes[second].priority) {
            nodes[first].right = self(self, nodes[first].right, second);
            pull(first); return first;
        }
        nodes[second].left = self(self, first, nodes[second].left);
        pull(second); return second;
    };
    int root = 0;
    for (ll value : values) {
        int node = nodes.size();
        nodes.push_back(Node{value, int(random() >> 1)});
        root = merge(merge, root, node);
    }
    for (auto [left, right] : reversals) {
        auto [first, rest] = split(split, root, left);
        auto [middle, last] = split(split, rest, right-left);
        if (middle) nodes[middle].reversed = !nodes[middle].reversed;
        root = merge(merge, merge(merge, first, middle), last);
    }
    vector<ll> answer;
    auto collect = [&](auto&& self, int node) -> void {
        if (!node) return;
        push(node);
        self(self, nodes[node].left); answer.push_back(nodes[node].value);
        self(self, nodes[node].right);
    };
    collect(collect, root);
    return answer;
}
\end{lstlisting}
\fi

% Разворачивает указанные $[left,right)$ в последовательности и возвращает
% итоговый массив. Локальные \texttt{split} и \texttt{merge} реализуют
% разрез по числу элементов и склейку; пустые отрезки допустимы.
% При случайных приоритетах ожидаемое время $O((n+q)\log(n+1))$,
% память $O(n)$; худшее время квадратичное. Фиксированный seed делает
% результаты воспроизводимыми, но не защищает от специально подобранного входа.
%
% Добавьте суммы и метки прибавления, вставку и удаление. Построение
% монотонным стеком уменьшает начальную работу до $O(n)$.

\iffalse
\section{Дерево Ли Чао для минимума прямых}
Запрос $(type,a,b)$ типа 0 добавляет прямую $y=ax+b$, типа 1 возвращает минимум при $x=a$; без прямых ответ \texttt{LLONG\_MAX}.
Все запросы известны заранее; $ax$ и $ax+b$ помещаются в \texttt{ll}, $q<\texttt{INT\_MAX}/4$.
Время $O(q\log(q+1))$, память $O(q)$.
\begin{lstlisting}[name=li_chao]
vector<ll> li_chao(const vector<tuple<int, ll, ll>>& queries) {
    vector<ll> coordinates;
    for (auto [type, a, b] : queries) {
        (void)b;
        if (type == 1) coordinates.push_back(a);
    }
    sort(coordinates.begin(), coordinates.end());
    coordinates.erase(unique(coordinates.begin(), coordinates.end()),
                      coordinates.end());
    if (coordinates.empty()) return {};
    int n = coordinates.size();
    vector<pair<ll, ll>> lines(4 * n);
    vector<bool> occupied(4 * n);
    auto evaluate = [&](pair<ll, ll> line, int position) {
        return line.first * coordinates[position] + line.second;
    };
    auto insert = [&](auto&& self, int node, int left, int right,
                      pair<ll, ll> line) -> void {
        if (!occupied[node]) {
            occupied[node] = true; lines[node] = line; return;
        }
        int middle = left + (right - left) / 2;
        bool betterLeft = evaluate(line, left) < evaluate(lines[node], left);
        bool betterMiddle = evaluate(line, middle) < evaluate(lines[node], middle);
        if (betterMiddle) swap(line, lines[node]);
        if (right - left == 1) return;
        if (betterLeft != betterMiddle) self(self, node*2, left, middle, line);
        else self(self, node*2+1, middle, right, line);
    };
    vector<ll> answer;
    for (auto [type, a, b] : queries) {
        if (type == 0) insert(insert, 1, 0, n, {a, b});
        else {
            int position = lower_bound(coordinates.begin(), coordinates.end(), a)
                         - coordinates.begin();
            ll best = LLONG_MAX;
            int node = 1, left = 0, right = n;
            while (true) {
                if (occupied[node]) best = min(best, evaluate(lines[node], position));
                if (right - left == 1) break;
                int middle = left + (right - left) / 2;
                if (position < middle) { node *= 2; right = middle; }
                else { node = node*2+1; left = middle; }
            }
            answer.push_back(best);
        }
    }
    return answer;
}
\end{lstlisting}
\fi

% Запрос $(type,a,b)$: тип 0 добавляет прямую $y=ax+b$, тип 1 находит
% минимум в $x=a$ (\texttt{b} игнорируется). Координаты запросов сжимаются
% заранее; ответ без прямых --- \texttt{LLONG\_MAX}.
% Каждое произведение $ax$ и значение $ax+b$ должно помещаться в \texttt{ll}.
% Время $O(q\log(q+1))$, память $O(q)$; $q<\texttt{INT\_MAX}/4$.
%
% Для онлайн-координат применяют динамическое дерево на заданном диапазоне.
% Можно добавлять прямую только на отрезке её применимости за $O(\log^2 n)$.

\iffalse
\section{Дерево отсортированных списков}
Для неизменяемого массива запрос $(left,right,value)$ возвращает число элементов $\le value$ на $[left,right)$.
Пустой массив допустим; $n<\texttt{INT\_MAX}/4$.
Построение и память $O(n\log(n+1))$, запрос $O(\log^2(n+1))$, ответы занимают $O(q)$ памяти.
\begin{lstlisting}[name=merge_sort_tree]
vector<int> merge_sort_tree(const vector<ll>& values,
        const vector<tuple<int, int, ll>>& queries) {
    int n = values.size();
    vector<vector<ll>> tree(4 * max(1, n));
    auto build = [&](auto&& self, int node, int left, int right) -> void {
        if (right - left == 1) { tree[node] = {values[left]}; return; }
        int middle = left + (right - left) / 2;
        self(self, node*2, left, middle); self(self, node*2+1, middle, right);
        merge(tree[node*2].begin(), tree[node*2].end(),
              tree[node*2+1].begin(), tree[node*2+1].end(),
              back_inserter(tree[node]));
    };
    auto count = [&](auto&& self, int node, int left, int right,
                     int queryLeft, int queryRight, ll value) -> int {
        if (queryRight <= left || right <= queryLeft) return 0;
        if (queryLeft <= left && right <= queryRight)
            return upper_bound(tree[node].begin(), tree[node].end(), value)
                 - tree[node].begin();
        int middle = left + (right - left) / 2;
        return self(self, node*2, left, middle, queryLeft, queryRight, value)
             + self(self, node*2+1, middle, right, queryLeft, queryRight, value);
    };
    if (n) build(build, 1, 0, n);
    vector<int> answer;
    for (auto [left, right, value] : queries)
        answer.push_back(n ? count(count, 1, 0, n, left, right, value) : 0);
    return answer;
}
\end{lstlisting}
% Запрос $(left,right,value)$ считает элементы $\le value$ на $[left,right)$.
% Массив неизменяемый. Построение и память $O(n\log(n+1))$,
% запрос $O(\log^2(n+1))$, ответы требуют $O(q)$ памяти.
% Пустой массив допустим; $n<\texttt{INT\_MAX}/4$.
%
% Дробное каскадирование ускоряет запросы до $O(\log n)$ с большей памятью.
% Для порядковых статистик можно использовать персистентное дерево частот.

\fi
\section{Монотонная очередь}
Для массива возвращает минимумы всех окон длины $width$; требуется $1\le width\le n$.
Время $O(n)$, дополнительная память $O(width)$ без ответов.
\begin{lstlisting}[name=monotonic_queue]
vector<ll> monotonic_queue(const vector<ll>& values, int width) {
    int n = values.size();
    assert(1 <= width && width <= n);
    deque<int> queue;
    vector<ll> answer;
    for (int i = 0; i < n; ++i) {
        while (!queue.empty() && queue.front() <= i-width) queue.pop_front();
        while (!queue.empty() && values[queue.back()] >= values[i]) queue.pop_back();
        queue.push_back(i);
        if (i >= width-1) answer.push_back(values[queue.front()]);
    }
    return answer;
}
\end{lstlisting}
% Возвращает минимумы всех окон длины $width$, где $1\le width\le n$.
% Каждый индекс добавляется и удаляется не более одного раза:
% время $O(n)$, дополнительная память $O(width)$ без ответов.
%
% Для максимума измените сравнение. Две очереди позволяют поддерживать
% разность максимума и минимума в окне с монотонно движущимися границами.

\iffalse
\section{Порядковые статистики вместо расширения PBDS}
Мультимножество изначально пусто; запрос $(type,value)$: 0 вставляет значение, 1 удаляет одно вхождение, если оно есть, 2 считает элементы строго меньше значения, 3 возвращает элемент с нулевым рангом \texttt{value}.
Ответы только для типов 2 и 3; ранг должен существовать, размеры поддеревьев помещаются в \texttt{int}.
Ожидаемое время $O(q\log(q+1))$, память $O(q)$; худшее время $O(q^2)$ и глубина рекурсии $O(q)$, фиксированный seed не защищает от специально подобранного входа.
\begin{lstlisting}[name=order_statistics]
vector<ll> order_statistics(const vector<pair<int, ll>>& queries) {
    struct Node {
        ll key; int priority;
        int left = 0, right = 0, count = 1, size = 1;
    };
    vector<Node> nodes(1, Node{0, 0});
    nodes[0].size = 0;
    mt19937 random(98123);
    auto pull = [&](int node) {
        nodes[node].size = nodes[node].count + nodes[nodes[node].left].size
                                           + nodes[nodes[node].right].size;
    };
    auto merge = [&](auto&& self, int first, int second) -> int {
        if (!first || !second) return first ? first : second;
        if (nodes[first].priority > nodes[second].priority) {
            nodes[first].right = self(self, nodes[first].right, second);
            pull(first); return first;
        }
        nodes[second].left = self(self, first, nodes[second].left);
        pull(second); return second;
    };
    auto change = [&](auto&& self, int node, ll key, int delta) -> int {
        if (!node) {
            if (delta < 0) return 0;
            int created = nodes.size();
            nodes.push_back(Node{key, int(random() >> 1)});
            return created;
        }
        if (key == nodes[node].key) {
            nodes[node].count += delta;
            if (!nodes[node].count)
                return merge(merge, nodes[node].left, nodes[node].right);
        } else {
            bool left = key < nodes[node].key;
            int child = self(self, left ? nodes[node].left : nodes[node].right,
                             key, delta);
            if (left) nodes[node].left = child;
            else nodes[node].right = child;
            if (child && nodes[child].priority > nodes[node].priority) {
                if (left) {
                    nodes[node].left = nodes[child].right;
                    nodes[child].right = node;
                } else {
                    nodes[node].right = nodes[child].left;
                    nodes[child].left = node;
                }
                pull(node); pull(child); return child;
            }
        }
        pull(node); return node;
    };
    int root = 0;
    vector<ll> answer;
    for (auto [type, value] : queries) {
        if (type < 2) root = change(change, root, value, type == 0 ? 1 : -1);
        else if (type == 2) {
            int node = root, count = 0;
            while (node) {
                if (value <= nodes[node].key) node = nodes[node].left;
                else {
                    count += nodes[nodes[node].left].size + nodes[node].count;
                    node = nodes[node].right;
                }
            }
            answer.push_back(count);
        } else {
            assert(0 <= value && value < nodes[root].size);
            int node = root;
            while (true) {
                int leftSize = nodes[nodes[node].left].size;
                if (value < leftSize) node = nodes[node].left;
                else if (value < leftSize + nodes[node].count) {
                    answer.push_back(nodes[node].key); break;
                } else {
                    value -= leftSize + nodes[node].count;
                    node = nodes[node].right;
                }
            }
        }
    }
    return answer;
}
\end{lstlisting}
% Стандартный treap без дополнительных заголовков заменяет
% \texttt{order\_of\_key} и \texttt{find\_by\_order} из PBDS.
% Мультимножество исходно пусто. Запрос $(type,value)$: 0 вставляет,
% 1 удаляет одно вхождение (если есть), 2 возвращает число элементов строго
% меньше \texttt{value}, 3 возвращает элемент с нулевым рангом \texttt{value}.
% Ранг должен существовать; ответы только для типов 2 и 3.
% Ожидаемое время $O(q\log(q+1))$, память $O(q)$, худшее время $O(q^2)$.
% Число элементов и сумма размеров помещаются в \texttt{int}.
% Фиксированный seed не даёт гарантии против специально подобранного входа.
%
% Если допустим GNU PBDS, нужны отдельные обязательные заголовки
% \path{ext/pb_ds/assoc_container.hpp} и
% \path{ext/pb_ds/tree_policy.hpp}; они не входят в
% \texttt{bits/stdc++.h}. Повторы там представляют парами из ключа и уникального
% номера. При известных координатах проще использовать дерево Фенвика.

\fi
\section{Тяжёло-лёгкая декомпозиция дерева}
Даны непустое неориентированное дерево с корнем 0 и значения вершин; запрос $(type,u,v,value)$ типа 0 присваивает \texttt{value} вершине $u$, типа 1 возвращает сумму на пути $u\ldots v$ включительно.
Требуется $n<\texttt{INT\_MAX}$; разности значений и все суммы помещаются в \texttt{ll}, нужен стек глубины $O(n)$.
Построение $O(n\log n)$, изменение $O(\log n)$, запрос пути $O(\log^2 n)$, память $O(n+q)$.
\begin{lstlisting}[name=hld]
vector<ll> hld(const vector<vector<int>>& graph, vector<ll> values,
        const vector<tuple<int, int, int, ll>>& queries) {
    int n = graph.size();
    vector<int> parent(n, -1), depth(n), size(n, 1), heavy(n, -1), head(n), position(n);
    auto sizes = [&](auto&& self, int vertex) -> void {
        for (int child : graph[vertex]) if (child != parent[vertex]) {
            parent[child] = vertex; depth[child] = depth[vertex] + 1;
            self(self, child); size[vertex] += size[child];
            if (heavy[vertex] == -1 || size[child] > size[heavy[vertex]])
                heavy[vertex] = child;
        }
    };
    sizes(sizes, 0);
    int timer = 0;
    auto decompose = [&](auto&& self, int vertex, int chain) -> void {
        head[vertex] = chain; position[vertex] = timer++;
        if (heavy[vertex] != -1) self(self, heavy[vertex], chain);
        for (int child : graph[vertex])
            if (child != parent[vertex] && child != heavy[vertex]) self(self, child, child);
    };
    decompose(decompose, 0, 0);
    vector<ll> tree(n + 1), answer;
    auto add = [&](int index, ll delta) {
        for (++index; index <= n;) {
            tree[index] += delta;
            int step = index & -index;
            if (step > n-index) break;
            index += step;
        }
    };
    auto prefix = [&](int right) {
        ll sum = 0;
        for (; right; right -= right & -right) sum += tree[right];
        return sum;
    };
    for (int i = 0; i < n; ++i) add(position[i], values[i]);
    for (auto [type, u, v, value] : queries) {
        if (type == 0) {
            add(position[u], value-values[u]); values[u] = value;
        } else {
            ll sum = 0;
            while (head[u] != head[v]) {
                if (depth[head[u]] < depth[head[v]]) swap(u, v);
                sum += prefix(position[u]+1) - prefix(position[head[u]]);
                u = parent[head[u]];
            }
            if (depth[u] > depth[v]) swap(u, v);
            sum += prefix(position[v]+1) - prefix(position[u]);
            answer.push_back(sum);
        }
    }
    return answer;
}
\end{lstlisting}
% Запрос $(type,u,v,value)$: тип 0 присваивает значение вершине $u$
% ($v$ игнорируется), тип 1 возвращает сумму на пути $u\ldots v$,
% включая обе вершины (\texttt{value} игнорируется).
% Начальные значения заданы массивом длины $n$.
% Построение $O(n\log n)$, изменение $O(\log n)$, путь $O(\log^2 n)$,
% память $O(n+q)$. $n<\texttt{INT\_MAX}$; разность нового и старого значений,
% суммы дерева Фенвика и все промежуточные суммы помещаются в \texttt{ll}.
%
% Для обновления всего пути замените Фенвик на ленивое дерево отрезков.
% Для некоммутативных агрегатов нужно учитывать направление каждого куска пути.

\section{Центроидная декомпозиция}
Для неориентированного дерева возвращает родителей в дереве центроидов; у корня родитель $-1$, пустой вход даёт пустой вектор.
Время $O(n\log n)$, память $O(n)$, высота результата $O(\log n)$; рекурсивным обходам нужен стек глубины $O(n)$.
\begin{lstlisting}[name=centroid_decomposition]
vector<int> centroid_decomposition(const vector<vector<int>>& graph) {
    int n = graph.size();
    vector<int> size(n), parent(n, -1);
    vector<bool> removed(n);
    auto measure = [&](auto&& self, int vertex, int previous) -> void {
        size[vertex] = 1;
        for (int child : graph[vertex]) if (child != previous && !removed[child]) {
            self(self, child, vertex); size[vertex] += size[child];
        }
    };
    auto locate = [&](auto&& self, int vertex, int previous, int total) -> int {
        for (int child : graph[vertex])
            if (child != previous && !removed[child] && size[child] > total/2)
                return self(self, child, vertex, total);
        return vertex;
    };
    auto build = [&](auto&& self, int vertex, int previous) -> void {
        measure(measure, vertex, -1);
        int center = locate(locate, vertex, -1, size[vertex]);
        parent[center] = previous; removed[center] = true;
        for (int child : graph[center]) if (!removed[child]) self(self, child, center);
    };
    if (n) build(build, 0, -1);
    return parent;
}
\end{lstlisting}
% Возвращает родителей в дереве центроидов; у корня родитель $-1$.
% Каждый центроид разбивает текущую компоненту на части размера не больше
% половины, поэтому высота результата $O(\log n)$.
% Время $O(n\log n)$, память $O(n)$. Пустой граф возвращает пустой вектор.
%
% Расстояния до предков-центроидов позволяют отвечать о ближайшей отмеченной
% вершине. Для подсчёта пар по расстоянию нужно вычитать вклад одной ветви.

\section{Виртуальное дерево}
Для дерева с корнем 0 и наборов важных вершин строит минимальные деревья, замкнутые по LCA.
Каждое ребро содержит предка, потомка и длину пути между ними.
Номера вершин исходные, рёбра направлены от предка к потомку, длина --- число исходных рёбер; повторы удаляются, пустой и одноэлементный наборы дают пустой список рёбер.
Предобработка $O(n\log n)$, набор размера $k$ --- $O(k\log(k+1)+k\log n)$, память $O(n\log n+\sum k)$; нужен стек глубины $O(n)$.
\begin{lstlisting}[name=virtual_tree]
vector<vector<tuple<int, int, int>>> virtual_tree(
        const vector<vector<int>>& graph, const vector<vector<int>>& groups) {
    int n = graph.size(), levels = 1, timer = 0;
    while ((1LL << levels) <= n) ++levels;
    vector<vector<int>> up(levels, vector<int>(n));
    vector<int> tin(n), tout(n), depth(n);
    auto dfs = [&](auto&& self, int vertex, int parent) -> void {
        tin[vertex] = timer++;
        up[0][vertex] = parent;
        for (int level = 1; level < levels; ++level)
            up[level][vertex] = up[level-1][up[level-1][vertex]];
        for (int child : graph[vertex]) if (child != parent) {
            depth[child] = depth[vertex]+1; self(self, child, vertex);
        }
        tout[vertex] = timer;
    };
    dfs(dfs, 0, 0);
    auto ancestor = [&](int a, int b) { return tin[a] <= tin[b] && tin[b] < tout[a]; };
    auto lca = [&](int a, int b) {
        if (ancestor(a, b)) return a;
        for (int level = levels-1; level >= 0; --level)
            if (!ancestor(up[level][a], b)) a = up[level][a];
        return up[0][a];
    };
    auto byEntry = [&](int a, int b) { return tin[a] < tin[b]; };
    vector<vector<tuple<int, int, int>>> answer;
    for (auto vertices : groups) {
        sort(vertices.begin(), vertices.end(), byEntry);
        vertices.erase(unique(vertices.begin(), vertices.end()), vertices.end());
        int count = vertices.size();
        for (int i = 1; i < count; ++i) vertices.push_back(lca(vertices[i-1], vertices[i]));
        sort(vertices.begin(), vertices.end(), byEntry);
        vertices.erase(unique(vertices.begin(), vertices.end()), vertices.end());
        vector<int> stack;
        vector<tuple<int, int, int>> edges;
        for (int vertex : vertices) {
            while (!stack.empty() && !ancestor(stack.back(), vertex)) stack.pop_back();
            if (!stack.empty()) edges.emplace_back(stack.back(), vertex,
                                                   depth[vertex]-depth[stack.back()]);
            stack.push_back(vertex);
        }
        answer.push_back(move(edges));
    }
    return answer;
}
\end{lstlisting}
% Для каждого набора важных вершин возвращает рёбра $(parent,child,length)$
% минимального замкнутого по LCA дерева. Вершины сохраняют исходные номера;
% рёбра направлены от предка к потомку, длина --- число исходных рёбер.
% Повторы во входном наборе удаляются; пустой и одноэлементный наборы дают
% пустой список рёбер. Корень непустого результата --- единственная вершина
% без входящего ребра; для одноэлементного набора это его единственная вершина.
% Предобработка $O(n\log n)$, набор размера $k$ обрабатывается за
% $O(k\log(k+1)+k\log n)$, память $O(n\log n+\sum k)$.
%
% Предобработку можно вынести наружу для потоков наборов. На виртуальном
% дереве удобно считать попарные расстояния и делать ДП по важным вершинам.

\iffalse
\section{Перекорневое динамическое программирование: суммы расстояний}
Для неориентированного невзвешенного дерева возвращает суммы расстояний от каждой вершины до остальных; пустой вход допустим.
Время и память $O(n)$; все суммы помещаются в \texttt{ll}, нужен стек глубины $O(n)$.
\begin{lstlisting}[name=reroot_distances]
vector<ll> reroot_distances(const vector<vector<int>>& graph) {
    int n = graph.size();
    if (!n) return {};
    vector<int> size(n, 1);
    vector<ll> answer(n);
    auto down = [&](auto&& self, int vertex, int parent) -> void {
        for (int child : graph[vertex]) if (child != parent) {
            self(self, child, vertex);
            size[vertex] += size[child];
            answer[vertex] += answer[child] + size[child];
        }
    };
    auto up = [&](auto&& self, int vertex, int parent) -> void {
        for (int child : graph[vertex]) if (child != parent) {
            answer[child] = answer[vertex] + n - 2LL*size[child];
            self(self, child, vertex);
        }
    };
    down(down, 0, -1); up(up, 0, -1);
    return answer;
}
\end{lstlisting}
% Для каждой вершины возвращает сумму расстояний до всех остальных в
% невзвешенном дереве. При переносе корня через ребро сумма меняется на
% $n-2\,size[child]$. Время и память $O(n)$; пустой граф допустим.
% Суммы и промежуточные результаты помещаются в \texttt{ll}.
%
% Для весов рёбер изменение умножается на вес. Для произвольной ДП
% используют префиксы и суффиксы вкладов соседей.

\fi
\iffalse
\section{Метод мешка: различные цвета в поддереве}
Для неориентированного дерева с корнем 0 и массива цветов длины $n$ возвращает число различных цветов в каждом поддереве.
Цвета имеют тип \texttt{ll}; пустые дерево и массив допустимы одновременно.
Время $O(n\log n)$, память $O(n)$, нужен стек глубины $O(n)$.
\begin{lstlisting}[name=dsu_on_tree]
vector<int> dsu_on_tree(const vector<vector<int>>& graph, const vector<ll>& colors) {
    int n = graph.size();
    if (!n) return {};
    vector<ll> coordinates = colors;
    sort(coordinates.begin(), coordinates.end());
    coordinates.erase(unique(coordinates.begin(), coordinates.end()), coordinates.end());
    vector<int> color(n), size(n, 1), heavy(n, -1), tin(n), tout(n), order;
    for (int i = 0; i < n; ++i)
        color[i] = lower_bound(coordinates.begin(), coordinates.end(), colors[i])
                 - coordinates.begin();
    auto prepare = [&](auto&& self, int vertex, int parent) -> void {
        tin[vertex] = order.size(); order.push_back(vertex);
        for (int child : graph[vertex]) if (child != parent) {
            self(self, child, vertex); size[vertex] += size[child];
            if (heavy[vertex] == -1 || size[child] > size[heavy[vertex]]) heavy[vertex] = child;
        }
        tout[vertex] = order.size();
    };
    prepare(prepare, 0, -1);
    vector<int> frequency(coordinates.size()), answer(n);
    int distinct = 0;
    auto change = [&](int left, int right, int delta) {
        for (int i = left; i < right; ++i) {
            int& count = frequency[color[order[i]]];
            if (count == 0) ++distinct;
            count += delta;
            if (count == 0) --distinct;
        }
    };
    auto solve = [&](auto&& self, int vertex, int parent, bool keep) -> void {
        for (int child : graph[vertex])
            if (child != parent && child != heavy[vertex]) self(self, child, vertex, false);
        if (heavy[vertex] != -1) self(self, heavy[vertex], vertex, true);
        for (int child : graph[vertex])
            if (child != parent && child != heavy[vertex]) change(tin[child], tout[child], 1);
        change(tin[vertex], tin[vertex]+1, 1);
        answer[vertex] = distinct;
        if (!keep) change(tin[vertex], tout[vertex], -1);
    };
    solve(solve, 0, -1, true);
    return answer;
}
\end{lstlisting}
% DSU on tree возвращает число различных цветов в поддереве каждой вершины.
% Цвета --- произвольные \texttt{ll}, предварительно сжимаются.
% Содержимое тяжёлого ребёнка сохраняется, лёгкие поддеревья добавляются
% по эйлеровым интервалам. Время $O(n\log n)$, память $O(n)$.
% Пустые граф и массив цветов допустимы одновременно.
%
% Частоты позволяют поддерживать число цветов с заданной частотой или
% наиболее частые цвета. В отличие от DSU связности, объединения здесь нет.

\fi
\iffalse
\section{Слияние меньшего множества в большее}
Для неориентированного дерева с корнем 0 и массива цветов длины $n$ возвращает число различных цветов в каждом поддереве, сливая меньшие множества в большие.
Пустой вход допустим; время $O(n\log^2(n+1))$, память $O(n)$, нужен стек глубины $O(n)$.
\begin{lstlisting}[name=small_to_large]
vector<int> small_to_large(const vector<vector<int>>& graph,
        const vector<ll>& colors) {
    int n = graph.size();
    vector<set<ll>> bags(n);
    vector<int> answer(n);
    auto dfs = [&](auto&& self, int vertex, int parent) -> void {
        bags[vertex].insert(colors[vertex]);
        for (int child : graph[vertex]) if (child != parent) {
            self(self, child, vertex);
            if (bags[vertex].size() < bags[child].size()) bags[vertex].swap(bags[child]);
            bags[vertex].insert(bags[child].begin(), bags[child].end());
            bags[child].clear();
        }
        answer[vertex] = bags[vertex].size();
    };
    if (n) dfs(dfs, 0, -1);
    return answer;
}
\end{lstlisting}
% Альтернативно считает различные цвета каждого поддерева, сливая
% \texttt{set} детей. Всегда переносится меньшее множество.
% Время $O(n\log^2(n+1))$, память $O(n)$; массив цветов имеет длину $n$.
% Пустой граф допустим. Дубликаты не мешают оценке: исчезающие элементы
% оплачивают одно слияние, а выжившие перемещения --- удвоением размера
% либо исчезновением сопоставимых элементов.
%
% Словарь частот заменяет множество для более богатой статистики.
% Хеш-таблица часто быстрее, но даёт только ожидаемые оценки и требует
% аккуратного выбора хеша на враждебном входе.

\fi
\section{Эйлеров порядок и интервалы поддеревьев}
Для неориентированного дерева с корнем 0 возвращает $(tin,tout,order)$: поддерево $v$ занимает $[tin[v],tout[v])$, \texttt{order[tin[v]]=v}.
Это порядок входов, не полный эйлеров обход для LCA; пустой вход допустим.
Время и память $O(n)$, рекурсии нет.
\begin{lstlisting}[name=euler_subtree]
tuple<vector<int>, vector<int>, vector<int>> euler_subtree(
        const vector<vector<int>>& graph) {
    int n = graph.size();
    vector<int> tin(n), tout(n), order;
    if (!n) return {tin, tout, order};
    vector<array<int, 3>> stack{{0, -1, 0}};
    while (!stack.empty()) {
        auto [vertex, parent, exit] = stack.back(); stack.pop_back();
        if (exit) { tout[vertex] = order.size(); continue; }
        tin[vertex] = order.size(); order.push_back(vertex);
        stack.push_back({vertex, parent, 1});
        for (auto it = graph[vertex].rbegin(); it != graph[vertex].rend(); ++it)
            if (*it != parent) stack.push_back({*it, vertex, 0});
    }
    return {tin, tout, order};
}
\end{lstlisting}
% Возвращает $(tin,tout,order)$: поддерево вершины $v$ занимает
% $[tin[v],tout[v])$, а \texttt{order[tin[v]]=v}.
% Время и память $O(n)$, обход итеративный и не требует рекурсивного стека.
% Пустой граф допустим. Это порядок входов, а не полный эйлеров обход для LCA.
%
% Дерево Фенвика поверх порядка даёт суммы поддеревьев. Для динамического
% изменения топологии нужны другие структуры, например деревья эйлерова тура.

\iffalse
\section{Изоморфизм неукоренённых деревьев}
Для двух неориентированных деревьев без меток возвращает признак точного изоморфизма без вероятностных хешей; два пустых дерева изоморфны.
Время $O(n\log(n+1))$, память $O(n)$, нужен стек глубины $O(n)$.
\begin{lstlisting}[name=tree_isomorphism]
bool tree_isomorphism(const vector<vector<int>>& first,
        const vector<vector<int>>& second) {
    if (first.size() != second.size()) return false;
    if (first.empty()) return true;
    auto centers = [&](const vector<vector<int>>& graph) {
        int n = graph.size(), remaining = n;
        vector<int> degree(n), leaves;
        for (int i = 0; i < n; ++i) {
            degree[i] = graph[i].size();
            if (degree[i] <= 1) leaves.push_back(i);
        }
        while (remaining > 2) {
            remaining -= leaves.size();
            vector<int> next;
            for (int vertex : leaves)
                for (int child : graph[vertex]) if (--degree[child] == 1) next.push_back(child);
            leaves = move(next);
        }
        return leaves;
    };
    map<vector<int>, int> identifiers;
    auto encode = [&](auto&& self, const vector<vector<int>>& graph,
                      int vertex, int parent) -> int {
        vector<int> children;
        for (int child : graph[vertex]) if (child != parent)
            children.push_back(self(self, graph, child, vertex));
        sort(children.begin(), children.end());
        int next = identifiers.size();
        return identifiers.emplace(move(children), next).first->second;
    };
    auto a = centers(first), b = centers(second);
    if (a.size() != b.size()) return false;
    int signature = encode(encode, first, a[0], -1);
    for (int root : b) if (signature == encode(encode, second, root, -1)) return true;
    return false;
}
\end{lstlisting}
% Проверяет точный изоморфизм без меток и без вероятностных хешей.
% Канонические номера присваиваются отсортированным спискам номеров детей
% в общем для обоих деревьев словаре. Проверяются один или два центра.
% Оба пустых дерева изоморфны. Время $O(n\log(n+1))$, память $O(n)$:
% суммарная длина ключей словаря линейна, а сравнения списков учитывают
% их длины; центров не больше двух. Рекурсия может иметь глубину $O(n)$.
%
% Для деревьев с метками включите метку вершины в ключ. Для восстановления
% биекции сопоставляйте детей с одинаковыми каноническими номерами.

\fi
\iffalse
\section{Центры и центроиды дерева}
Для неориентированного невзвешенного дерева возвращает пару отсортированных списков: центры минимизируют максимальное расстояние, центроиды оставляют после удаления компоненты размера $\le n/2$.
Каждый список содержит не больше двух вершин; пустой вход допустим.
Время и память $O(n)$, рекурсии нет.
\begin{lstlisting}[name=tree_centers]
pair<vector<int>, vector<int>> tree_centers(const vector<vector<int>>& graph) {
    int n = graph.size();
    if (!n) return {{}, {}};
    vector<int> parent(n, -1), order{0}, size(n, 1), degree(n), centers, centroids;
    for (int i = 0; i < n; ++i) {
        int vertex = order[i];
        for (int child : graph[vertex]) if (child != parent[vertex]) {
            parent[child] = vertex; order.push_back(child);
        }
        degree[i] = graph[i].size();
        if (degree[i] <= 1) centers.push_back(i);
    }
    for (int i = n-1; i >= 0; --i) {
        int vertex = order[i], largest = n-size[vertex];
        for (int child : graph[vertex]) if (parent[child] == vertex)
            largest = max(largest, size[child]);
        if (largest <= n/2) centroids.push_back(vertex);
        if (parent[vertex] != -1) size[parent[vertex]] += size[vertex];
    }
    int remaining = n;
    while (remaining > 2) {
        remaining -= centers.size();
        vector<int> next;
        for (int vertex : centers)
            for (int child : graph[vertex]) if (--degree[child] == 1) next.push_back(child);
        centers = move(next);
    }
    sort(centers.begin(), centers.end()); sort(centroids.begin(), centroids.end());
    return {centers, centroids};
}
\end{lstlisting}
% Возвращает пару отсортированных списков: центры минимизируют максимальное
% расстояние до других вершин, центроиды оставляют после удаления компоненты
% размера не больше $n/2$. Это разные понятия; каждого типа вершин не больше
% двух. Время и память $O(n)$, обходы итеративные; пустое дерево допустимо.
%
% Центры также являются серединой диаметра. Взвешенные расстояния требуют
% другого поиска центра, а взвешенные размеры --- адаптации критерия центроида.

\fi
